\documentclass[11pt,a4paper]{article}
\usepackage[utf8]{inputenc}
\usepackage[margin=1in]{geometry}
\usepackage{amsmath}
\usepackage{booktabs}
\usepackage{hyperref}
\usepackage{xcolor}
\usepackage{tcolorbox}
\usepackage{enumitem}

\hypersetup{colorlinks=true, linkcolor=blue!60!black, urlcolor=blue!60!black}
\pagestyle{empty}

\begin{document}

\begin{center}
{\Large\bfseries Why the Galleri Trial Missed Its Primary Endpoint}\\[0.3em]
{\large\itshape Diagnostic Delay, Not Test Failure}\\[0.5em]
{\normalsize February 2026}
\end{center}

\subsection*{The Finding}

The NHS-Galleri trial's primary endpoint---a reduction in combined
Stage~III+IV cancer diagnoses---was not met. Our analysis shows why.

Using a Monte Carlo simulation calibrated to NHS population data, we
estimate that the Galleri test correctly detected approximately
\textbf{84 cancers at Stage~I or~II} that were subsequently recorded
as Stage~III or~IV because the NHS diagnostic pathway took too long
to confirm them. These cancers progressed past the Stage~II/III
boundary during workup---imaging, biopsies, specialist referrals,
and MDT review---and were counted \emph{against} the very endpoint
the test was trying to improve.

\begin{tcolorbox}[colback=blue!3!white, colframe=blue!50!black,
  boxrule=0.5pt, arc=2pt, left=6pt, right=6pt, top=4pt, bottom=4pt]
The trial needed just \textbf{25 of these 84 cases}---roughly one in
three---to have been diagnosed promptly for the primary endpoint to
reach statistical significance ($p < 0.05$). The test's biological
performance is intact. The miss is an infrastructure problem, not an
assay problem.
\end{tcolorbox}

\subsection*{The Context}

The trial enrolled 142{,}000 participants across England and screened
for 12 cancer types using the Galleri multi-cancer blood test. The
intervention arm showed clear success: Stage~IV diagnoses fell by
approximately 30\%, and Stage~I/II detections rose by over 50\%. But
the composite Stage~III+IV endpoint---the primary measure---showed
only about 470 cases in the control arm versus 443 in the
intervention arm. A difference of just 26 cases, nowhere near
statistical significance.

The puzzle: how can a test that demonstrably finds cancers earlier
fail to reduce late-stage diagnoses?

\subsection*{The Mechanism: Stage Slip}

When the blood test flags a possible cancer, the patient enters the
NHS diagnostic pathway. This process takes time---our analysis
estimates a median of 92~days for standard-of-care workup. During
those weeks and months, the tumour continues to grow. A cancer that
was Stage~II when the test detected it may be Stage~III by the time
the diagnosis is finalised. We call this \textbf{stage slip}: the
test found the cancer early, but the system recorded it late.

Critically, this is a problem unique to the intervention arm. In the
control arm, patients with no screening test simply develop cancer
along its natural course until symptoms eventually trigger a
referral---typically when the disease is already advanced. That is
the baseline the trial measures against, not an infrastructure
failure. Stage slip only occurs when there \emph{was} an early
detection event that the diagnostic system failed to capitalise on.

\subsection*{The Evidence Behind the Numbers}

Our simulation is not a theoretical exercise. It is built on
observed data at every level:

\begin{itemize}[nosep, leftmargin=2em]
\item \textbf{Calibrated to NHS reality.} For each of the 12 cancer
  types, the model's control arm reproduces the actual
  stage-at-diagnosis distribution reported in NHS England statistics
  (NLCA, NHS Digital, Cancer Research~UK)---typically within 1--2
  percentage points. The baseline is an actuarial match for England.

\item \textbf{Cancer-specific detection rates} from the published
  CCGA~3 clinical validation study, the largest multi-cancer
  screening study to date.

\item \textbf{A 92-day diagnostic delay} constructed from five
  independent evidence streams: optimised US trial benchmarks
  (PATHFINDER~1 and~2), NHS waiting-time performance data, a
  published study of trial-induced diagnostic spillover across 9.6
  million NHS referrals, and GRAIL's own regulatory disclosure.

\item \textbf{200 independent simulation runs} producing statistical
  confidence intervals on every output.
\end{itemize}

\subsection*{Robustness}

The conclusion does not depend on the exact delay assumption. We
tested every value from 65 to 120~days:

\begin{center}\small
\begin{tabular}{rl}
\toprule
SoC Delay & Composite endpoint (RR, 95\% CI) \\
\midrule
65 days  & 0.911 \,[CI spans 1.0] \\
92 days  & 0.943 \,[CI spans 1.0] \\
120 days & 0.981 \,[CI spans 1.0] \\
\bottomrule
\end{tabular}
\end{center}

\noindent At every delay value tested, the composite endpoint remains
non-significant. A critic who argues the true delay is 75~days
rather than 92 has moved a parameter without changing the conclusion.

GRAIL's own press release---subject to SEC financial disclosure
standards---independently confirms the mechanism: ``there was a
higher than anticipated incidence of Stage~III cancers'' linked to
``time to diagnostic resolution,'' which ``appears to improve over
time as physicians gain experience.''

\subsection*{The Implication}

Our simulation shows that realistic NHS diagnostic delays are large
 enough to fully account for the primary endpoint miss. The number
  of cancers estimated to “slip” stages exceeds what was needed for
   the trial to meet its success threshold.
%The results are consistent with the test finding cancers earlier.
%%The Galleri test works.
% %It finds cancers earlier. 
 The NHS
diagnostic pathway---never designed for asymptomatic MCED referrals
and operating under significant capacity constraints during the trial
period---was too slow to translate early detection into early
\emph{diagnosis} for a material number of cases. The trial's 12-month
extension, combined with the learning-curve improvements GRAIL has
already documented, provides the window for this gap to close.

The endpoint was not missed because the test failed. It was missed
because the infrastructure was not yet fast enough to keep up with
what the test found.

\end{document}
